\documentclass[10pt]{article}
\usepackage[margin=0.9in]{geometry}
\usepackage{amsmath,amssymb,amsthm}
\usepackage[hidelinks]{hyperref}
\newtheorem{theorem}{Theorem}
\title{Moving counting jumps disprove a\\periodic-oscillation conjecture}
\author{Principia Math In Progress}
\date{18 September 2026}
\begin{document}
\maketitle
\begin{abstract}
Give each nonempty binary word of length \(n\) the size
\(n+\sqrt2\,2^{-n}\). Its generating functions satisfy the analytic
hypotheses of Bevan and Cond\'e's Conjecture 6, including a full
lattice of simple dominant poles. Nevertheless its cumulative count
has no asymptotic of the asserted periodic form. This is a
counterexample to MathDB 366299 and to the conjecture in the
published article.
\end{abstract}

\section*{Task 1: a counterexample to the original conjecture}
Conjecture 6 of \cite{BC}, represented by MathDB 366299, predicts
\[
 |\mathcal C_{\le x}|\sim c(x)\rho^{-x}x^{\alpha-1}
 \qquad (x\longrightarrow\infty)
\]
with bounded \(\omega\)-periodic \(c\), under two conditions:
the irrational generating function has the decomposition
\[
 f(z)=g(z)+\frac{h(z)}{(1-z/\rho)^\alpha},
 \quad 0<\rho<1,\quad \alpha\notin-\mathbb N,
\]
where \(g,h\) are analytic on the closed cut disk
\(\overline D(0,\rho)\setminus\mathbb R_{\le0}\); and the only
singularities of the Dirichlet generating function on its convergence
line are \(\log(1/\rho)+2\pi i\mathbb Z/\omega\).
The counting variable \(x\) is real, as in the source's definition.

\begin{theorem}
Let \(\mathcal C\) consist of the nonempty binary words, with each
word of length \(n\ge1\) assigned size
\[
 \lambda_n=n+a2^{-n},\qquad a=\sqrt2.
\]
This is an irrational combinatorial class satisfying all these
hypotheses with \(\rho=1/2\), \(\alpha=1\), \(h=1\), and
\(\omega=1\). There is no bounded \(1\)-periodic function \(c\)
such that \( |\mathcal C_{\le x}|\sim c(x)2^x\).
\end{theorem}

\begin{proof}
\emph{The class.}
There are \(2^n\) words of size \(\lambda_n\). Since
\(0<a2^{-n}<1\), we have \(n<\lambda_n<n+1\), so the sizes
increase strictly, have no finite accumulation point, and have
finite multiplicities. This is precisely the required notion of a
combinatorial class. It is irrational even in the source's stronger
sense of not having all sizes in any \(\eta\mathbb N\), \(\eta>0\).
Indeed, if \(\lambda_1/\lambda_2=r\in\mathbb Q\), then
\(1+a/2=r(2+a/4)\); comparing the rational and irrational parts
would give both \(r=1/2\) and \(r=2\).

\emph{The irrational generating function.}
Use the principal logarithm on \(\mathbb C\setminus\mathbb R_{\le0}\).
Then
\[
 f(z)=\sum_{n\ge1}2^n z^{n+a2^{-n}}.
\]
For fixed \(z\ne0\), the absolute value of its \(n\)th summand is
\((2|z|)^n |z|^{a2^{-n}}\), with the last factor tending to \(1\).
Thus the series converges for \(|z|<1/2\), and its terms fail to
tend to zero for \(|z|\ge1/2\). Its radius is exactly \(1/2\).

On \(|z|<1/2\) off the cut, subtracting a geometric series gives
\begin{align*}
 f(z)&=\frac{1}{1-2z}+g(z),\\
 g(z)&=-1+\sum_{n\ge1}2^n z^n
              \bigl(\exp(a2^{-n}\operatorname{Log}z)-1\bigr).
\end{align*}
We show that \(g\) is analytic on the larger open set
\(U=\{z:|z|<1\}\setminus\mathbb R_{\le0}\).
On any compact subset of \(U\), take \(|z|\le r<1\) and
\(|\operatorname{Log}z|\le L\). The elementary estimate
\(|e^w-1|\le |w|e^{|w|}\) bounds the \(n\)th summand by
\[
 2^n r^n a2^{-n}L e^{a2^{-n}L}
 \le aL e^{aL/2}r^n.
\]
The series therefore converges uniformly on compact subsets of
\(U\). Each summand is analytic there, hence so is \(g\).
In particular \(g\) is analytic on a neighborhood of every point
of the required closed cut disk. The choice \(h=1\), \(\alpha=1\)
meets the stated hypothesis, with \(h(1/2)=1\ne0\).

\emph{The entire line of dominant poles.}
For \(\operatorname{Re}s>\log2\), the Dirichlet series satisfies
\begin{align*}
 F(s)&=\sum_{n\ge1}2^n e^{-s(n+a2^{-n})}
      =\frac{2e^{-s}}{1-2e^{-s}}+H(s),\\
 H(s)&=\sum_{n\ge1}2^n e^{-ns}
                  \bigl(e^{-as2^{-n}}-1\bigr).
\end{align*}
Its abscissa is \(\log2\), by the same absolute-value test.
On a compact subset of \(\operatorname{Re}s>0\), take
\(\operatorname{Re}s\ge\sigma>0\) and \(|s|\le M\). Now the
\(n\)th summand of \(H\) has absolute value at most
\[
 aM e^{aM/2}e^{-n\sigma}.
\]
Consequently \(H\) is analytic throughout this larger half-plane.
The displayed expression continues \(F\) meromorphically across
its entire convergence line. Its singularities there are exactly
\[
 s_k=\log2+2\pi i k,\qquad k\in\mathbb Z.
\]
At each \(s_k\), the numerator \(2e^{-s}\) and the derivative of
the denominator \(1-2e^{-s}\) both equal \(1\). Every pole is
simple, has residue \(1\), and cannot be canceled by the analytic
function \(H\). This verifies the full pole condition with
\(\omega=1\), without substituting across a logarithmic branch.

\emph{The asserted periodic asymptotic is impossible.}
Write \(A(x)=|\mathcal C_{\le x}|\), and suppose that
\(A(x)/(c(x)2^x)\to1\) for a \(1\)-periodic \(c\).
Fix any \(r\in(0,1)\). For every sufficiently large integer \(N\),
\(a2^{-N}<r\), so
\[
 A(N+r)=\sum_{n=1}^N2^n=2^{N+1}-2.
\]
Since \(c(N+r)=c(r)\), the assumed asymptotic forces
\[
 c(r)=2^{1-r}\quad\hbox{for every }r\in(0,1).
\]
At integers it similarly forces \(c(0)=1\), because
\(A(N)=2^N-2\). These conclusions require no regularity of \(c\).

Set \(\delta_N=a2^{-N-1}\) and \(x_N=N+\delta_N\).
Here \(N<x_N<\lambda_N\), so \(A(x_N)=2^N-2\), whereas
periodicity and the already forced values on \((0,1)\) imply
\(c(x_N)=2^{1-\delta_N}\). Thus
\[
 \frac{A(x_N)}{c(x_N)2^{x_N}}
 =\frac{2^N-2}{2^{N+1}}
 =\frac12-2^{-N}\longrightarrow\frac12,
\]
contradicting the required limit \(1\).
\end{proof}

\subsection*{Scope and verification}
The failure comes from counting jumps whose phases move toward zero.
The limit along each fixed fractional part exists, but those limits
do not produce an asymptotic valid for all real \(x\to\infty\).
The example uses positive integral multiplicities, genuinely
irrational sizes, a nonvanishing singular coefficient, and identical
simple principal parts at every boundary pole. The source permits
general locally finite size sets; it does not require the size of a
word to be the sum of fixed letter weights.

The proof is self-contained. The accompanying verifier independently
counts with exact arithmetic in \(\mathbb Q(\sqrt2)\), checks the
local residue symbolically, and compares two convergent representations
of the generating functions numerically. The numerical checks are
supplemental; analyticity and the limiting contradiction follow from
the displayed estimates and identities for all \(n\).
The complete current preprint (29 pages) and journal article
(31 pages) were read. Conjecture 6 remains unchanged in the final.
The separate source and literature audits record the bounded
eligibility checks. Review was by the authoring agent; no independent
peer review or proof-assistant certification is claimed.

\section*{Task 2: fixed lattice phases}
Establish a corrected theorem for size sets contained in a fixed
finite union of translates of \(\omega\mathbb N\). Specify the
coefficient asymptotics needed to make the resulting periodic
approximation valid at every real counting threshold.

\section*{Task 3: smoothed counting functions}
Determine which averaging kernels convert the boundary pole data into
a periodic leading term. Quantify how the answer depends on the
averaging width, including widths that tend to zero.

\section*{Task 4: stability under perturbing sizes}
For multiplicities \(b^n\) at sizes \(n+\varepsilon_n\), classify
the perturbations that preserve a periodic cumulative asymptotic.
Describe the exceptional intervals near jumps when such an asymptotic
holds away from those intervals, and estimate their total length.

\section*{Task 5: additional analytic conditions}
Identify boundary growth or uniform Tauberian hypotheses that rule
out moving jumps. Pole positions and even their full principal parts
alone are insufficient, as this example shows.
Tasks 2--5 are proposed research directions and carry no extra credit.

\begin{thebibliography}{9}
\bibitem{BC}
D. Bevan and J. Cond\'e,
\emph{Introducing irrational enumeration: analytic combinatorics for
objects of irrational size},
Electronic Journal of Combinatorics \textbf{33}(3) (2026), P3.15,
Conjecture 6, p.~15.
\url{https://doi.org/10.37236/14759}.
Current preprint: arXiv:2412.14682v3 (22 December 2025),
Conjecture 6, p.~14.
\url{https://arxiv.org/abs/2412.14682v3}.
\end{thebibliography}
\end{document}
